\documentclass[10pt,a4paper]{amsart}
\usepackage[margin=19mm]{geometry}
\usepackage{amsmath,amssymb,mathtools}
\usepackage[scr=rsfs,cal=euler]{mathalfa}
\usepackage{xfrac}
\usepackage[unicode,hidelinks,bookmarks=false]{hyperref}
\newcommand{\ie}{i.e.\ }
\newcommand{\cf}{cf.\ }
\newcommand{\resp}{respectively\ }
\newcommand{\Subsection}{\subsection}
\providecommand{\nobreakdash}{\nobreak\hbox{-}}
\setlength{\emergencystretch}{2em}
\multlinegap0pt
\usepackage{amssymb}
\usepackage{float}
\usepackage{xcolor}
\usepackage{subfig}
\newtheorem{theorem}{Theorem}
\title{Herman Rings: Structure, Dynamics, and Open Problems}
\author{Gorachand Chakraborty, Subhasis Ghora, Tarakanta Nayak}
\date{Source: arXiv:2512.24118v1 (CC BY 4.0)}
\begin{document}
\maketitle

\noindent\emph{Source and licence.} Herman Rings: Structure, Dynamics, and Open Problems. Version arXiv:2512.24118v1 (CC BY 4.0), distributed by arXiv under the Creative Commons Attribution 4.0 International licence, http://creativecommons.org/licenses/by/4.0/, as recorded in the arXiv licence metadata for this exact version and shown on the abstract page https://arxiv.org/abs/2512.24118v1. Full licence notice: \texttt{licenses/arxiv-2512.24118-CC-BY-4.0.txt}.\par\smallskip

\noindent\textit{Editorial scope.} Counted unit: the article's theorem nohr (source0.tex lines 314--317). The article's complete original proof of it is delivered with it (source0.tex lines 318--320, one \texttt{proof} environment of 3 lines). No mathematics is written, repaired or re-derived: the statement and the proof are the article's own lines, in the article's own notation, and no retained line is substituted or re-flowed.  No further source-local context is needed: every cross-reference of every retained line resolves inside the two delivered spans.  Nothing was omitted from any retained span, so the package carries no omission record.  No retained line contains a citation, so the wrapper carries no bibliography and no bibliography entry.  No retained line contains a figure environment, an image-inclusion command or drawing code, so no image is delivered and the package declares no asset dependency.  Editorial formatting only, in addition: an AMS \texttt{amsart} preamble that restates the article's own declarations and macros used by the retained lines, namely the environments \texttt{theorem} on separate counters, together with the standard packages those lines need.  The retained environments print in this document as Theorem 1 in order of appearance, whereas the article prints the same items with the numbering of its own full document.  The complete native source of the article as distributed in the arXiv e-print for this exact version is bundled unchanged as \texttt{sources/191854/source0.tex}.
\begin{theorem}
\label{nohr}
    If $f$ is a polynomial (or an entire function), then $f$ has no Herman rings.
\end{theorem}

\begin{proof}
    Let $f$ be a polynomial (or an entire function), and if possible, let $H$ be a $p$-periodic Herman ring of $f$. Also, let $\gamma$ be an $f^p$-invariant Jordan curve of $H$. Since $f^p(\gamma)=\gamma$, either $f^p(B(\gamma))=B(\gamma)$ or $f^p(B(\gamma))=\widehat{\mathbb{C}}\setminus (B(\gamma)\cup \gamma)$. The map  $f^p$ is a polynomial (or an entire function); hence, $B(\gamma)$ does not contain any pole of $f^p$. Thus $f^p(B(\gamma))=B(\gamma)$ giving that $f^{np}(B(\gamma))=B(\gamma)$ for $n\in \mathbb{N}$. This is not possible as $B(\gamma)\cap J(f)\neq \emptyset$ giving that $\cup_{n\in \mathbb{N}}f^{np}(B(\gamma))$ is dense in $\widehat{\mathbb{C}}$. This concludes that $f$ does not have any Herman ring. 
\end{proof}

\end{document}
